\section{Theoretical Applications}\label{sec:theoretical-applications}

\subsection{A Minority Principle}
Let us consider two-point opinion distributions 
supported on the extremes $\{0,1\}$ where one mass is smaller than the other, for example 
$\pi=((10,90),(0,1)).$ 
One might expect that the share of the total $(\alpha,\beta)$-effort of each mass required to reach an optimal consensus $y^*$ for $\pi$ should depend on the parameters $\alpha$
and $\beta$. In particular, it would seem natural that, by varying these
parameters, one could shift the burden of adjustment between the masses:
for some values, the larger mass might bear most of the effort, while
for others, the smaller mass might.

From the definition of $\EC$ the effort of the smaller mass to $y^*$ is $10^\alpha {y^*}^\beta$ while for the larger mass we obtain $90^\alpha (1-y^*)^\beta$. If $\beta=1$ the larger mass does not move at all, since from Cor.~\ref{cor:mec-special-cases}, we have $y^* = 1$. Thus, the smaller mass bears the whole effort. The situation is more subtle when $\beta>1$. In this case, the optimal
consensus point lies in the interior $(0,1)$, and thus both masses contribute to the
effort.  

Nevertheless, and somewhat surprisingly, it is actually the smaller mass that \emph{always} makes the greater effort to reach the optimal consensus, regardless of $\alpha$ and $\beta$. This is established below for any two-point distribution \(\pi\) with support \(\{l,r\}\), where \(l<r\).

First the following lemma confirms the intuition that the optimal consensus point lies closer to the larger mass.

\begin{restatable}{lemma}{yoptClosest}\label{lem:yopt-closest-to-largest-general}
Assume \(\beta>1\). Let \(0\le l<r\le 1\), and let \(\pi\in\Dset\) satisfy \(\pi(l)=A\) and \(\pi(r)=B\), with \(A,B>0\). Let \(y^*\) be such that \(\MEC(\pi)=\EC[\alpha,\beta][\pi](y^*)\). Then:
\begin{enumerate}
\item if \(y^*>\frac{l+r}{2}\), then \(B>A\);
\item if \(y^*<\frac{l+r}{2}\), then \(B<A\);
\item if \(y^*=\frac{l+r}{2}\), then \(B=A\).
\end{enumerate}
\end{restatable}

Having located the side of the optimal consensus point relative to the two masses, we now compare their respective contributions to the total effort.

\begin{restatable}{lemma}{greatestContribution}\label{lem:greatest-contribution-from-smallest-general}
Under the assumptions of Lem.~\ref{lem:yopt-closest-to-largest-general}, we have:
\begin{enumerate}
\item if \(y^*>\frac{l+r}{2}\), then \(A^\alpha (y^*-l)^\beta > B^\alpha (r-y^*)^\beta\);
\item if \(y^*<\frac{l+r}{2}\), then \(A^\alpha (y^*-l)^\beta < B^\alpha (r-y^*)^\beta\);
\item if \(y^*=\frac{l+r}{2}\), then \(A^\alpha (y^*-l)^\beta = B^\alpha (r-y^*)^\beta\).
\end{enumerate}
\end{restatable}

As an immediate consequence of the above lemmas and Cor.~\ref{cor:mec-special-cases}, we conclude that indeed it is the minority that always bears the greater effort.

\begin{restatable}[Minority Principle for Two-Point Distributions]{theorem}{minorityPrinciple}
\label{thm:minority-principle-general}
Let \(0\le l<r\le 1\), and let \(\pi\in\Dset\) satisfy \(\pi(l)=A\) and \(\pi(r)=B\), with \(A,B>0\) and \(A\neq B\). Let \(y^*\) be such that \(\MEC(\pi)=\EC[\alpha,\beta][\pi](y^*)\). Then:
\[
A<B \;\Longrightarrow\; A^\alpha (y^*-l)^\beta > B^\alpha (r-y^*)^\beta,
\]
and
\[
A>B \;\Longrightarrow\; A^\alpha (y^*-l)^\beta < B^\alpha (r-y^*)^\beta.
\]
\end{restatable}

\subsection{A Tipping Point Method}\label{sec:tipping-point} 

In Ex.~\ref{ex:group-identity}, we illustrated the role of group identity using a distribution $\pi$ and showing that transferring a small portion of mass at 0.75 toward the extreme right may in fact decrease polarization. Such a shift can reduce the effective weight of the original mass and thereby lower the overall level of polarization.

Nevertheless, if $\pi'$ is obtained from $\pi$ in Ex.~\ref{ex:group-identity} by transferring the entire mass
at $0.75$ to the extreme right, polarization does increase. Indeed, we have 
\(
\MEC[2,2](\pi)<\MEC[2,2](\pi') \approx 461.5 .
\)
This naturally suggests the following question: what is the minimal amount of
mass (\emph{critical mass}) that must move from the moderate right (0.75) to the extreme right (1.0) for
polarization to increase? In other words, at what point does such a shift begin
to raise polarization?

We now illustrate a simple method to solve the above tipping-point question for a family of two-point distributions with group identity and alienation parameters $\alpha=\beta=2.$ The phenomenon illustrates a broader fact: under MEC, polarization is not monotone in movement away from the optimal consensus point. Whether a partial transfer of mass away from this point raises or lowers polarization depends on how much mass is involved, and the Tipping Point method makes this dependence explicit and computable.


\begin{example}[A tipping point method for mass shift]
\label{ex:tipping-mass-shift-general}
Let \(m,n>0\) and \(0\le l<r<r'\le 1\). Consider the opinion distribution
\(
\pi=((m,n),(l,r)),
\)
and, for \(k\in(0,n]\), let
\(
\pi_{r:k\rightarrow r'}=((m,n-k,k),(l,r,r'))
\)
be the distribution obtained by moving mass \(k\) from \(r\) to the more extreme point \(r'\).

To study how polarization changes under this transfer, define
\[
\Delta(k)\defsymbol \MEC[2,2](\pi_{r:k\rightarrow r'})-\MEC[2,2](\pi).
\]

Set \(d\defsymbol r-l\), \(e\defsymbol r'-r\), and \(C\defsymbol m^2+n^2\). By Proposition~\ref{prop:mec-general-mass-shift-formulas}, we have
\[
\Delta(k)=\frac{k\,R_{m,n,d,e}(k)}{Q_{m,n}(k)},
\]
where
\[
R_{m,n,d,e}(k)
=
e^2Ck^3-2e^2nCk^2+\bigl(2d^2m^4+2de\,m^2C+e^2C^2\bigr)k-2d^2m^4n,
\]
and
\[
Q_{m,n}(k)=C(2k^2-2nk+C).
\]

Since \(Q_{m,n}(k)>0\) for all \(k\in[0,n]\), the sign of \(\Delta(k)\) for \(k>0\) is determined by \(R_{m,n,d,e}(k)\). Moreover,
\[
R_{m,n,d,e}(0)=-2d^2m^4n<0,
\qquad
R_{m,n,d,e}(n)=e(2d+e)m^2n(m^2+n^2)>0.
\]
By Proposition~\ref{prop:unique-root-tipping-general}, the polynomial \(R_{m,n,d,e}\) has a unique root \(k^\star\in(0,n)\). Therefore
\[
\MEC[2,2](\pi_{r:k^\star\rightarrow r'})=\MEC[2,2](\pi),
\]
and
\[
\MEC[2,2](\pi_{r:k\rightarrow r'})<\MEC[2,2](\pi)
\quad\text{for }0<k<k^\star,
\]
while
\[
\MEC[2,2](\pi_{r:k\rightarrow r'})>\MEC[2,2](\pi)
\quad\text{for }k^\star<k\le n.
\]

Thus \(k^\star\) acts as a tipping point: transferring a small amount of mass from \(r\) to \(r'\) initially reduces polarization, but once the transferred mass exceeds \(k^\star\), the polarization becomes strictly larger than that of the original distribution. \qed
\end{example}

\begin{remark}
By the Reflection Lemma (Lem.~\ref{lem:mec-reflection}), the analogous case in which a mass \(k\), with \(0<k\le m\), is moved leftward from \(l\) to a more extreme point \(l'<l\), reduces to the rightward case by reflection about \(c=(l'+r)/2\). Indeed, let
\[
\pi_{l:k\to l'}=((k,m-k,n),(l',l,r)).
\]
If \(r''\defsymbol l'+r-l\), then \(l'<r''<r\), \(\pi^{\mathrm{ref}(c)}=((n,m),(l',r''))\), and \((\pi_{l:k\to l'})^{\mathrm{ref}(c)}=((n,m-k,k),(l',r'',r))\). Thus the result follows from the rightward case.
\end{remark}

We now show an application of the tipping point method to the distribution in 
Ex.~\ref{ex:group-identity}.

\begin{example}
Let
\[
\pi=((40,60),(0.25,0.75)),
\]
so that
\[
\pi_{0.75:k\rightarrow1}=((40,60-k,k),(0.25,0.75,1)).
\]
Applying the analysis of Ex.~\ref{ex:tipping-mass-shift-general} with
\(m=40\), \(n=60\), \(l=0.25\), \(r=0.75\), and \(r'=1\), we obtain
\(d=\frac12\) and \(e=\frac14\). Thus the change in polarization is
determined by the cubic polynomial
\[
R_{40,60,\frac12,\frac14}(k)
=
325k^3-39000k^2+5050000k-76800000.
\]
This polynomial has a unique real root in \((0,60)\) 
\[
k^\star\approx 17.156.
\]
Therefore
\[
\MEC[2,2](\pi_{0.75:k\rightarrow1})<\MEC[2,2](\pi)
\quad\text{for }0<k<k^\star,
\]
while
\[
\MEC[2,2](\pi_{0.75:k\rightarrow1})>\MEC[2,2](\pi)
\quad\text{for }k^\star<k\le 60.
\]

Thus the polarization becomes strictly larger than that of \(\pi\) once the
transferred mass exceeds approximately \(17.156\).
\end{example}

We conclude by showing that the tipping point method above extends beyond the two-point setting. For any opinion distribution and any mass transferred from its maximal support point toward the extreme right (\(1\)), the resulting change in \(\MEC[2,2]\) is again a rational function whose sign is governed by a cubic polynomial.

\begin{restatable}[Mass transfer to the extreme]{proposition}{massTransfer}
\label{prop:mec22-mass-transfer}
Let
\(
\pi=((p_1,\dots,p_r),(x_1,\dots,x_r))
\)
be an opinion distribution with
\(
0\le x_1<\cdots<x_r<1
\)
and \(p_i>0\). For \(k\in[0,p_r]\), let \(\pi^{(k)}=\pi_{x_r:k\to 1}\) be the distribution obtained by moving mass \(k\) from the maximal support point \(x_r\) to \(1\). Then there exist polynomials \(P_\pi(k)\) and \(Q_\pi(k)\) such that
\[
\MEC[2,2](\pi^{(k)})-\MEC[2,2](\pi)
=
\frac{k\,P_\pi(k)}{Q_\pi(k)}
\qquad\text{for all }k\in[0,p_r],
\]
where \(P_\pi(k)\) is a cubic polynomial and \(Q_\pi(k)>0\) for all \(k\in[0,p_r]\). Consequently, for every \(k>0\) the sign of the change in polarization is determined by the cubic polynomial \(P_\pi(k)\).
\end{restatable}
